\section{Самоорганизующиеся карты Кохонена}
\label{sec:som}

\noindent\hspace{1.25cm}\emph{Раздел используется в ЛР~№4.}

\subsection{Конкурентное обучение}

В задаче кластеризации нет целевых выходов, поэтому сеть должна извлечь структуру из самих входных
векторов. Простейший механизм~--- \emph{конкурентное обучение}. Слой содержит $M$ нейронов с векторами
весов $W_1, \ldots, W_M$ той же размерности, что и вход. На вход $\mathbf{p}$ откликается один нейрон~---
\emph{победитель} (best matching unit, BMU), вектор весов которого ближе всего к входу:
\begin{equation}
\label{eq:bmu}
c = \argmin_{i} \lVert W_i - \mathbf{p}\rVert .
\end{equation}
По правилу <<победитель получает всё>> (winner-take-all, WTA) корректируются только его веса:
$W_c \leftarrow W_c + \eta(\mathbf{p} - W_c)$. Нейрон-победитель сдвигается к входу на долю $\eta$
расстояния, и нейрон, выигрывающий на некоторой группе точек, в итоге оказывается в её центре. По сути
это стохастический вариант алгоритма $k$-средних ($k$-means), минимизирующий \emph{квантизационную
ошибку}~--- среднее расстояние от точек до ближайшего нейрона. Недостаток WTA~--- \emph{<<мёртвые>>
нейроны}: нейрон, изначально оказавшийся далеко от данных, никогда не побеждает и не обучается.

\subsection{Архитектура и алгоритм самоорганизующейся карты}

Самоорганизующаяся карта (self-organizing map, SOM), предложенная Т.~Кохоненом~\cite{kohonen2008},
дополняет конкурентное обучение \emph{топологией}. Нейроны расположены в узлах решётки, и у каждого
нейрона $i$, кроме вектора весов $W_i$ в пространстве данных, есть фиксированное положение $\mathbf{r}_i$
на решётке. Биологический прототип~--- топографические карты коры головного мозга, где соседние участки
реагируют на близкие стимулы (соматотопическая, ретинотопическая, тонотопическая организация).

При обучении сдвигаются веса не только победителя, но и его соседей по решётке:
\begin{equation}
\label{eq:som}
W_i(t + 1) = W_i(t) + \eta(t)\,h_{ci}(t)\,\bigl(\mathbf{p} - W_i(t)\bigr),
\end{equation}
где $h_{ci}(t)$~--- \emph{функция окрестности}, убывающая с расстоянием между нейронами $c$ и $i$
\emph{на решётке}. Наиболее распространена гауссова окрестность
\begin{equation}
h_{ci}(t) = \exp\left(-\frac{\lVert \mathbf{r}_c - \mathbf{r}_i\rVert^2}{2\sigma^2(t)}\right).
\end{equation}
Скорость обучения $\eta(t)$ и радиус окрестности $\sigma(t)$ уменьшаются в ходе обучения, например
экспоненциально или линейно: $\sigma$ от половины диаметра решётки до долей шага, $\eta$ от
$0,5$ до $0,01$. В начале, пока окрестность широкая, вся карта движется согласованно и
\emph{упорядочивается}: соседние на решётке нейроны занимают соседние области пространства данных. В
конце окрестность сужается до одного нейрона, и веса точно \emph{подстраиваются} под плотность данных.
Благодаря соседству в обучении участвуют и нейроны, далёкие от данных, поэтому мёртвых нейронов
становится значительно меньше.

\emph{Топология решётки}: прямоугольная (\code{gridtop}), у внутреннего нейрона
4 ближайших соседа; гексагональная (\code{hextop}), нечётные ряды сдвинуты на половину шага, у
внутреннего нейрона 6 соседей на одинаковом расстоянии. \emph{Расстояние на решётке} измеряют евклидовой
метрикой между координатами узлов или числом связей по кратчайшему пути (\code{linkdist},
рисунок~\ref{fig:som-topology}). Гексагональная решётка изотропнее прямоугольной: окрестность ближе к кругу,
и карта изгибается равномернее.

\begin{figure}[!htb]
\centering
\includegraphics[width=0.8\textwidth]{\figdir/som_topology.png}
\caption{Прямоугольная и гексагональная решётки 5$\times$5: числа~--- расстояние \code{linkdist}
от центрального нейрона}
\label{fig:som-topology}
\end{figure}

\emph{Алгоритм обучения SOM}:
\begin{enumerate}
\item задать решётку, инициализировать веса случайно в области данных (или по главным компонентам);
\item для каждой эпохи и каждого входа $\mathbf{p}$ в случайном порядке: найти победителя
(\ref{eq:bmu}), вычислить $h_{ci}$ для всех нейронов, обновить веса по формуле (\ref{eq:som});
\item после эпохи уменьшить $\eta$ и $\sigma$; повторять до заданного числа эпох;
\item отнести каждую точку к кластеру своего нейрона-победителя.
\end{enumerate}

На практике применяют и более простое правило с двумя фазами (так работает, например, \code{learnsom}
в MATLAB): в \emph{фазе упорядочивания} скорость падает с $0,9$ до $0,02$, а радиус окрестности~--- от
максимального расстояния на решётке до фиксированного значения (по умолчанию 1), в \emph{фазе
подстройки} оба параметра постоянны; соседи в пределах радиуса получают шаг вдвое меньше победителя.
Если окрестность не сужается до одного нейрона, на маленькой карте все нейроны остаются соседями до конца
обучения, продолжают тянуть друг друга, и центроиды кластеров стягиваются к центру данных
(рисунок~\ref{fig:som-rules}). Поэтому для точного нахождения центров кластеров радиус к концу обучения
уменьшают до нуля или используют гауссову окрестность с убывающим $\sigma$.

\begin{figure}[!htb]
\centering
\includegraphics[width=\textwidth]{\labfig{4}{07_rules_2x2.png}}
\caption{Карта 2$\times$2 при разных правилах обучения: несужающаяся окрестность стягивает
центроиды (ЛР~№4)}
\label{fig:som-rules}
\end{figure}

\subsection{Оценка качества карты}

Поскольку истинные метки кластеров в реальных задачах неизвестны, качество карты оценивают
внутренними показателями:
\begin{itemize}
\item \emph{квантизационная ошибка} $QE = \frac{1}{N}\sum_{k}\lVert \mathbf{p}_k - W_{c(\mathbf{p}_k)}
\rVert$~--- точность представления данных нейронами; убывает с ростом числа нейронов;
\item \emph{топографическая ошибка} $TE$~--- доля точек, для которых первый и второй ближайшие нейроны
не являются соседями на решётке; показывает, насколько карта сохраняет топологию;
\item \emph{карта попаданий} (hits)~--- число точек, для которых нейрон является победителем; нейроны без
попаданий лежат между кластерами или вне данных;
\item \emph{U-матрица}~--- для каждого нейрона среднее расстояние между его весами и весами соседей;
большие значения образуют <<хребты>> на границах кластеров.
\end{itemize}
Если истинное разбиение известно, совпадение с ним измеряют индексом Рэнда (ARI, 1~--- полное совпадение).

Когда число нейронов равно числу кластеров, SOM эквивалентна $k$-средним. \emph{Избыточная} карта
растягивается по данным: нейроны распределяются внутри облаков пропорционально их плотности, а
несколько нейронов остаются между облаками, удерживаемые связями решётки; на U-матрице они образуют
границы кластеров.

\subsection{Реализация на Python}
\label{sec:py-som}

Шаг обучения карты (\ref{eq:som}) для одного входа \code{p}: \code{W}~--- веса $M \times 2$, \code{R}~---
координаты узлов решётки $M \times 2$ (для \code{hextop} нечётные ряды сдвинуты на $0,5$ по горизонтали
и сжаты по вертикали в $\sqrt 3/2$ раз):
\begin{lstlisting}
c = np.argmin(((W - p) ** 2).sum(axis=1))                         # нейрон-победитель
h = np.exp(-((R - R[c]) ** 2).sum(axis=1) / (2 * sigma ** 2))     # окрестность на решётке
W += eta * h[:, None] * (p - W)                                   # сдвиг победителя и соседей
\end{lstlisting}
Для сравнения служит библиотека MiniSom~\cite{vettigli2018}; она же вычисляет показатели качества:
\begin{lstlisting}
from minisom import MiniSom
som = MiniSom(5, 5, 2, sigma=2.0, learning_rate=0.5, topology='hexagonal',
              neighborhood_function='gaussian', random_seed=3)
som.random_weights_init(X)
som.train(X, 200 * len(X), random_order=True)        # 200 эпох
qe, te = som.quantization_error(X), som.topographic_error(X)
U, hits = som.distance_map(), som.activation_response(X)   # U-матрица, попадания
\end{lstlisting}
Карту с числом нейронов, равным числу кластеров, полезно сравнить с \code{sklearn.cluster.KMeans}, а
совпадение разбиений измерить функцией \code{sklearn.metrics.adjusted\_rand\_score}.
